\documentclass[10pt,a4paper]{amsart}
\usepackage[margin=19mm]{geometry}
\usepackage{amsmath,amssymb,mathtools}
\usepackage[scr=rsfs,cal=euler]{mathalfa}
\usepackage{xfrac}
\usepackage[unicode,hidelinks,bookmarks=false]{hyperref}
\newcommand{\ie}{i.e.\ }
\newcommand{\cf}{cf.\ }
\newcommand{\resp}{respectively\ }
\newcommand{\Subsection}{\subsection}
\providecommand{\nobreakdash}{\nobreak\hbox{-}}
\setlength{\emergencystretch}{2em}
\multlinegap0pt
\usepackage{amssymb}
\newtheorem{lemma}{Lemma}
\title{(Claw, $\vec C_3$)-free Digraphs with Unbounded Dichromatic Number}
\author{Guillaume Aubian, Luis Kuffner}
\date{Source: arXiv:2602.08736v1 (CC BY 4.0)}
\begin{document}
\maketitle

\noindent\emph{Source and licence.} (Claw, $\vec C_3$)-free Digraphs with Unbounded Dichromatic Number. Version arXiv:2602.08736v1 (CC BY 4.0), distributed by arXiv under the Creative Commons Attribution 4.0 International licence, http://creativecommons.org/licenses/by/4.0/, as recorded in the arXiv licence metadata for this exact version and shown on the abstract page https://arxiv.org/abs/2602.08736v1. Full licence notice: \texttt{licenses/arxiv-2602.08736-CC-BY-4.0.txt}.\par\smallskip

\noindent\textit{Editorial scope.} Counted unit: the article's lemma lemma (source0.tex lines 58--60). The article's complete original proof of it is delivered with it (source0.tex lines 61--66, one \texttt{proof} environment of 6 lines). No mathematics is written, repaired or re-derived: the statement and the proof are the article's own lines, in the article's own notation, and no retained line is substituted or re-flowed.  No further source-local context is needed: every cross-reference of every retained line resolves inside the two delivered spans.  Nothing was omitted from any retained span, so the package carries no omission record.  No retained line contains a citation, so the wrapper carries no bibliography and no bibliography entry.  No retained line contains a figure environment, an image-inclusion command or drawing code, so no image is delivered and the package declares no asset dependency.  Editorial formatting only, in addition: an AMS \texttt{amsart} preamble that restates the article's own declarations and macros used by the retained lines, namely the environments \texttt{lemma} on separate counters, together with the standard packages those lines need.  The retained environments print in this document as Lemma 1 in order of appearance, whereas the article prints the same items with the numbering of its own full document.  The complete native source of the article as distributed in the arXiv e-print for this exact version is bundled unchanged as \texttt{sources/194268/source0.tex}.
\begin{lemma}
When $|a - c| = |b - d|$, vertices $(a,b),\ (a,d),\ (c,d),\ (c,b)$ induce a directed $4$-cycle.
\end{lemma}

\begin{proof}
Let $i=\min \{j \mid |a - c|_{j}  = 1\}$. Then $i = \min\{j \mid a_j \ne c_j\} = \min\{j \mid b_j \ne d_j\}$.
Up to permuting $a$ and $c$ and/or permuting $b$ and $d$, we may assume $a_i = b_i = 1$ and $c_i = d_i = 0$.
Since $b_i = a_i$, we have $(a,b) \to (a,d)$. Since $d_i \ne a_i$, we have $(a,d) \to (c,d)$. Since $d_i = c_i$, we have $(c,d) \to (c,b)$. Since $b_i \ne c_i$, we have $(c,b) \to (a,b)$.
Thus the four vertices form the directed cycle $(a,b) \to (a,d) \to (c,d) \to (c,b) \to (a,b)$.
\end{proof}

\end{document}
